\documentclass[12pt,a4paper]{article}

% --- PAQUETES DE IDIOMA Y FORMATO ---
\usepackage[utf8]{inputenc}
\usepackage[spanish,es-tabla]{babel}
\usepackage{geometry}
\geometry{top=2.5cm, bottom=2.5cm, left=3cm, right=2.5cm}
\usepackage{graphicx}
\usepackage{booktabs}
\usepackage{tabularx}
\usepackage{longtable}
\usepackage{array}
\usepackage{colortbl}
\usepackage{xcolor}
\usepackage{hyperref}
\usepackage{enumitem}
\usepackage{fancyhdr}
\usepackage{listings}

% --- COLORES INSTITUCIONALES UNSAAC ---
\definecolor{unsaacBlue}{RGB}{0, 51, 102}
\definecolor{unsaacGold}{RGB}{204, 153, 0}
\definecolor{codeBg}{RGB}{245, 247, 250}
\definecolor{codeBorder}{RGB}{218, 224, 233}
\definecolor{darkgray}{RGB}{60, 60, 60}

% --- CONFIGURACIÓN DE LISTINGS PARA CÓDIGO DART ---
\lstdefinelanguage{Dart}{
  morekeywords={abstract,as,assert,async,await,break,case,catch,class,const,continue,default,deferred,do,dynamic,else,enum,export,extends,extension,external,factory,false,final,finally,for,Function,get,hide,if,implements,import,in,interface,is,late,library,mixin,new,null,of,on,operator,part,rethrow,return,set,show,static,super,switch,sync,this,throw,true,try,typedef,var,void,while,with,yield},
  sensitive=true,
  morecomment=[l]{//},
  morecomment=[s]{/*}{*/},
  morestring=[b]',
  morestring=[b]"
}

\lstset{
  language=Dart,
  backgroundcolor=\color{codeBg},
  basicstyle=\ttfamily\footnotesize\color{darkgray},
  keywordstyle=\bfseries\color{unsaacBlue},
  commentstyle=\itshape\color{gray},
  stringstyle=\color{unsaacGold!80!black},
  breaklines=true,
  frame=single,
  rulecolor=\color{codeBorder},
  numbers=left,
  numberstyle=\tiny\color{gray},
  numbersep=8pt,
  showstringspaces=false,
  tabsize=2
}

% --- ENCABEZADOS Y PIES DE PÁGINA ---
\pagestyle{fancy}
\fancyhf{}
\fancyhead[L]{\small\color{gray}UNSAAC - Desarrollo de Aplicaciones Móviles (IF616AIN)}
\fancyhead[R]{\small\color{gray}Informe N° 01}
\fancyfoot[C]{\thepage}
\renewcommand{\headrulewidth}{0.4pt}

\hypersetup{
    colorlinks=true,
    linkcolor=unsaacBlue,
    citecolor=unsaacBlue,
    urlcolor=unsaacBlue
}

\begin{document}

% --- PORTADA INSTITUCIONAL ---
\begin{titlepage}
    \centering
    \vfill
    {\large\bfseries UNIVERSIDAD NACIONAL DE SAN ANTONIO ABAD DEL CUSCO} \\[0.4cm]
    {\large FACULTAD DE INGENIERÍA ELÉCTRICA, ELECTRÓNICA, INFORMÁTICA Y MECÁNICA} \\[0.4cm]
    {\large ESCUELA PROFESIONAL DE INGENIERÍA INFORMÁTICA Y DE SISTEMAS} \\[1.2cm]
    
    \rule{0.4\textwidth}{0.6pt} \\[0.3cm]
    {\bfseries [ LOGO INSTITUCIONAL UNSAAC ]} \\[0.3cm]
    \rule{0.4\textwidth}{0.6pt} \\[1.2cm]
    
    {\Large\bfseries ASIGNATURA: DESARROLLO DE APLICACIONES MÓVILES} \\[0.3cm]
    {\large\bfseries CÓDIGO: IF616AIN \quad | \quad SEMESTRE: 2026-1} \\[0.8cm]
    
    {\Large\bfseries INFORME DE LABORATORIO N° 01} \\[0.4cm]
    {\large\bfseries PRODUCTO 1: INSTALACIÓN DEL ENTORNO DE DESARROLLO Y APLICACIÓN BÁSICA ``HOLA MUNDO MÓVIL''} \\[1.2cm]
    
    \textbf{DOCENTE RESPONSABLE:} \\
    Ing. CCACYAHUILLCA-BEJAR-HANS HARLEY \\[0.8cm]
    
    \textbf{ESTUDIANTE / DESARROLLADOR:} \\
    SUPA CUSIPAUCAR, Yeferson \quad (Código: 220553) \\[1.2cm]
    
    {\large Cusco - Perú} \\
    {\large 2026}
    \vfill
\end{titlepage}

\newpage
\tableofcontents
\newpage

% --- 1. INTRODUCCIÓN ---
\section{Introducción y Objetivos}

\subsection{Contexto de la Asignatura}
La asignatura \textbf{Desarrollo de Aplicaciones Móviles (IF616AIN)} de la Universidad Nacional de San Antonio Abad del Cusco (UNSAAC) tiene por propósito brindar los fundamentos teóricos y prácticos indispensables para la creación de software móvil moderno, robusto y eficiente. 

De acuerdo con la Primera Unidad Didáctica (\textit{Fundamentos del desarrollo móvil}), el primer hito de aprendizaje corresponde a la instalación de las herramientas del ecosistema de desarrollo y la construcción del \textbf{Producto 1: Aplicación básica ``Hola Mundo móvil''}.

\subsection{Objetivos de la Práctica}
\begin{itemize}[leftmargin=*]
    \item \textbf{Objetivo General:} Configurar de forma integral el entorno de desarrollo móvil y construir una primera aplicación móvil en Flutter y Dart que evidencie el ciclo de vida y la reactividad de componentes.
    \item \textbf{Objetivos Específicos:}
    \begin{enumerate}
        \item Instalar y configurar el SDK de Flutter (canal stable) junto con el Java Development Kit (JDK 17) y el Android SDK (plataforma 36 y build-tools).
        \item Integrar el entorno de desarrollo en Visual Studio Code y Android Studio con sus respectivas extensiones y herramientas de depuración.
        \item Comprender la arquitectura de directorios de un proyecto Flutter y la función de sus archivos esenciales (\texttt{pubspec.yaml}, \texttt{lib/main.dart}, \texttt{android/}).
        \item Implementar widgets de tipo \texttt{StatelessWidget} y \texttt{StatefulWidget}, aplicando el concepto de reactividad mediante \texttt{setState()}.
        \item Realizar pruebas automatizadas y ejecutar la aplicación en entornos de prueba y dispositivos móviles reales.
    \end{enumerate}
\end{itemize}

% --- 2. FUNDAMENTOS TEÓRICOS ---
\section{Fundamentos del Desarrollo Móvil}

\subsection{Tipos de Aplicaciones Móviles}
\begin{itemize}[leftmargin=*]
    \item \textbf{Nativas:} Desarrolladas en los lenguajes oficiales de cada plataforma (Kotlin/Java para Android, Swift para iOS). Ofrecen el máximo rendimiento y acceso a APIs de bajo nivel, pero requieren código duplicado.
    \item \textbf{Híbridas / Web:} Aplicaciones empaquetadas basadas en tecnologías web (HTML, CSS, JavaScript) dentro de un contenedor WebView. Tienen menor rendimiento al depender de un motor web intermediario.
    \item \textbf{Multiplataforma Compiladas (Flutter):} Framework de Google que utiliza el lenguaje Dart y compila directamente a código binario de máquina (ARM/x86) utilizando el motor gráfico propio (Impeller / Skia). Permite un rendimiento idéntico al nativo con una única base de código para Android, iOS, Web y Escritorio.
\end{itemize}

\subsection{Arquitectura y Ciclo de Vida de una Aplicación Flutter}
En Flutter, todo elemento visual es un \textbf{Widget}. La interfaz de usuario es reactiva: se describe en función del estado actual (\textit{UI = f(state)}).
\begin{itemize}[leftmargin=*]
    \item \texttt{StatelessWidget}: Widget inmutable cuya estructura no depende de variables mutables a lo largo del tiempo.
    \item \texttt{StatefulWidget}: Widget asociado a una clase de estado (\texttt{State}) que persiste en memoria y permite reconstruir la interfaz ante eventos de usuario mediante la invocación de \texttt{setState()}.
\end{itemize}

% --- 3. CONFIGURACIÓN DEL ENTORNO ---
\section{Instalación y Configuración del Entorno de Desarrollo}

En concordancia con las herramientas digitales estipuladas en el sílabo, se procedió con la instalación y configuración técnica de los siguientes componentes:

\subsection{Herramientas Instaladas}
\begin{table}[h!]
\centering
\small
\begin{tabularx}{\textwidth}{lXp{5.5cm}}
\toprule
\textbf{Herramienta} & \textbf{Versión Instalada} & \textbf{Ruta / Configuración} \\
\midrule
\textbf{Flutter SDK} & 3.47.4 (Canal stable) & \texttt{C:\textbackslash src\textbackslash flutter\textbackslash bin} \\
\textbf{Dart SDK} & 3.13.3 & Integrado con Flutter \\
\textbf{Java JDK} & Eclipse Temurin 17.0.20.1 & \texttt{JAVA\_HOME} configurado \\
\textbf{Android Studio} & Versión Oficial Google 2026.1.4 & Con accesos directos configurados \\
\textbf{Android SDK} & Plataforma 36, Build-Tools 36.0.0 & \texttt{ANDROID\_HOME} configurado \\
\textbf{Platform-Tools} & adb 1.0.41 (v37.0.1) & Para depuración en terminal \\
\textbf{VS Code} & Extensiones Flutter y Dart 3.142.0 & Integrado para Hot Reload \\
\textbf{Postman} & 12.28.0 & Para consumo de APIs REST \\
\textbf{Figma Desktop} & 126.9.7 & Para prototipado UI/UX \\
\textbf{Firebase CLI} & 15.30.1 & Para servicios cloud (Unidad II) \\
\bottomrule
\end{tabularx}
\caption{Resumen de herramientas del entorno de desarrollo móvil.}
\end{table}

\subsection{Auditoría del Entorno (\texttt{flutter doctor})}
Para garantizar que todos los componentes se encuentren enlazados sin conflictos, se ejecutó la herramienta de diagnóstico \texttt{flutter doctor -v}, obteniendo un resultado limpio y sin observaciones:

\begin{lstlisting}[language=bash, numbers=none]
[√] Flutter (Channel stable, 3.47.4, on Microsoft Windows, locale es-PE)
[√] Windows Version (Windows 11 or higher)
[√] Android toolchain - develop for Android devices (Android SDK version 36.0.0)
    • Android SDK at C:\Users\yefer\AppData\Local\Android\Sdk
    • Platform android-36, build-tools 36.0.0
    • ANDROID_HOME = C:\Users\yefer\AppData\Local\Android\Sdk
    • All Android licenses accepted.
[√] Chrome - develop for the web
[√] Visual Studio - develop Windows apps (Visual Studio Community 2022)
[√] Connected device (3 available)
[√] Network resources

• No issues found!
\end{lstlisting}

% --- 4. ESTRUCTURA Y DESARROLLO ---
\section{Desarrollo del Producto 1: Hola Mundo Móvil}

\subsection{Estructura de Carpetas del Proyecto}
El proyecto generado mediante \texttt{flutter create} organiza su código de acuerdo con el estándar de la industria:
\begin{itemize}[leftmargin=*]
    \item \texttt{lib/}: Directorio principal donde reside el código fuente en Dart.
    \item \texttt{lib/main.dart}: Archivo de entrada ejecutable de la aplicación.
    \item \texttt{pubspec.yaml}: Archivo de manifiesto donde se declaran las dependencias externas, versiones del SDK, tipografías y recursos multimedia.
    \item \texttt{android/}: Contenedor del proyecto nativo Android (Gradle, configuración de paquete \texttt{pe.edu.unsaac.hola\_mundo}, permisos en \texttt{AndroidManifest.xml}).
    \item \texttt{test/}: Módulo de pruebas de regresión, integración y pruebas de widgets.
\end{itemize}

\subsection{Implementación de la Aplicación (\texttt{lib/main.dart})}
A continuación se presenta el código implementado, el cual contextualiza la información de la UNSAAC, la asignatura y demuestra la reactividad mediante un carrusel dinámico de mensajes y contador interactivo:

\begin{lstlisting}
import 'package:flutter/material.dart';

void main() {
  runApp(const UnsaacMobileApp());
}

/// Widget raiz inmutable: define el tema visual Material 3
class UnsaacMobileApp extends StatelessWidget {
  const UnsaacMobileApp({super.key});

  @override
  Widget build(BuildContext context) {
    return MaterialApp(
      title: 'UNSAAC - Desarrollo Movil',
      debugShowCheckedModeBanner: false,
      theme: ThemeData(
        useMaterial3: true,
        colorScheme: ColorScheme.fromSeed(
          seedColor: const Color(0xFF003366), // Azul institucional
          brightness: Brightness.light,
        ),
      ),
      home: const HolaMundoScreen(),
    );
  }
}

/// Pantalla principal reactiva (StatefulWidget)
class HolaMundoScreen extends StatefulWidget {
  const HolaMundoScreen({super.key});

  @override
  State<HolaMundoScreen> createState() => _HolaMundoScreenState();
}

class _HolaMundoScreenState extends State<HolaMundoScreen> {
  int _contadorSaludos = 0;
  final List<String> _mensajes = [
    'Hola Mundo movil! 📱',
    'Bienvenido a Desarrollo de Aplicaciones Moviles',
    'UNSAAC - Ingenieria Informatica y de Sistemas',
    'Construido con Flutter y Dart en tiempo record 🚀',
  ];
  int _indiceMensaje = 0;

  void _cambiarMensaje() {
    // Notifica al framework que el estado cambio y redibuja la interfaz
    setState(() {
      _contadorSaludos++;
      _indiceMensaje = (_indiceMensaje + 1) % _mensajes.length;
    });
  }

  @override
  Widget build(BuildContext context) {
    final theme = Theme.of(context);
    return Scaffold(
      appBar: AppBar(
        title: const Text('IF616AIN - UNSAAC', style: TextStyle(fontWeight: FontWeight.bold)),
        centerTitle: true,
        backgroundColor: theme.colorScheme.primaryContainer,
      ),
      body: SafeArea(
        child: SingleChildScrollView(
          padding: const EdgeInsets.all(20.0),
          child: Column(
            crossAxisAlignment: CrossAxisAlignment.stretch,
            children: [
              Card(
                elevation: 4,
                child: Padding(
                  padding: const EdgeInsets.all(24.0),
                  child: Column(
                    children: [
                      Icon(Icons.smartphone_rounded, size: 72, color: theme.colorScheme.primary),
                      const SizedBox(height: 16),
                      Text(
                        _mensajes[_indiceMensaje],
                        textAlign: TextAlign.center,
                        style: theme.textTheme.headlineSmall?.copyWith(fontWeight: FontWeight.bold),
                      ),
                      const SizedBox(height: 12),
                      const Text('Producto 1: Aplicacion basica "Hola Mundo movil"'),
                    ],
                  ),
                ),
              ),
              const SizedBox(height: 20),
              FilledButton.icon(
                onPressed: _cambiarMensaje,
                icon: const Icon(Icons.touch_app_rounded),
                label: const Text('Presioname para cambiar el saludo'),
              ),
            ],
          ),
        ),
      ),
    );
  }
}
\end{lstlisting}

% --- 5. PRUEBAS Y VALIDACIÓN ---
\section{Pruebas Automatizadas y Validación}

\subsection{Pruebas de Widgets (\texttt{test/widget\_test.dart})}
Se implementó una prueba automatizada para verificar el correcto renderizado inicial del saludo institucional y la respuesta del estado ante la pulsación del botón:

\begin{lstlisting}[language=bash, numbers=none]
$ flutter test
00:00 +0: loading test/widget_test.dart
00:00 +0: Prueba de humo para UnsaacMobileApp
00:00 +1: All tests passed!
\end{lstlisting}

\subsection{Análisis Estático de Código (\texttt{flutter analyze})}
Se ejecutó el analizador estático de Dart para garantizar el cumplimiento de las buenas prácticas y guías de estilo oficiales:
\begin{lstlisting}[language=bash, numbers=none]
$ flutter analyze
Analyzing Desarrollo_SoftwareII...
No issues found! (ran in 7.5s)
\end{lstlisting}

% --- 6. CONCLUSIONES ---
\section{Conclusiones}
\begin{itemize}[leftmargin=*]
    \item Se completó con éxito la instalación y configuración del entorno de desarrollo móvil para la asignatura \textbf{IF616AIN}, validando la total interoperabilidad entre Flutter 3.47, Java 17, Android Studio y el Android SDK 36.
    \item La aplicación desarrollada cumple con los requerimientos del \textbf{Producto 1}, permitiendo comprobar el flujo de compilación, la estructura modular de directorios y la reactividad de componentes mediante widgets \texttt{Stateless} y \texttt{Stateful}.
    \item El mecanismo de \textit{Hot Reload} de Flutter optimiza sustancialmente los tiempos de desarrollo de interfaz frente al ciclo tradicional de compilación nativa en Android.
\end{itemize}

% --- 7. REFERENCIAS ---
\section{Referencias Bibliográficas}
\begin{itemize}[leftmargin=*]
    \item Martín, F. (2022). \textit{Desarrollo de aplicaciones móviles con Flutter y Dart}. Anaya Multimedia.
    \item Deitel, P., \& Deitel, H. (2021). \textit{Android for Programmers: An App-Driven Approach}. Prentice Hall.
    \item Documentación Oficial de Flutter. \url{https://flutter.dev/docs}
    \item Google Developers Portal. \url{https://developer.android.com}
\end{itemize}

\end{document}
